% Discussion, limitations, and ethics.
% Author: fangtoast
% Email: fangtoast@foxmail.com
% Copyright (c) 2026 fangtoast. All rights reserved.

\section{Discussion, Limitations, and Ethics}

The experiments suggest two complementary conclusions. First, Fourier-space
global coupling and joint full-band prediction provide a useful inductive bias:
BISON is more accurate than substantially larger dense models. Second, forward
accuracy is necessary but not sufficient for inverse design. The 0.118-point
surrogate--replay gap is small relative to the final error, yet individual
targets still fail more severely and the independent query model can disagree.
Retaining the MPB boundary is therefore part of the method, not an optional
post-processing detail.

The results also clarify why the method should be described as a system rather
than as a single forward benchmark. Full-band accuracy makes BISON useful for
ranking objectives that depend on several bands and wave vectors at once. The
query operator probes whether this score is consistent under a second output
factorization, while MPB decides whether the selected geometry actually
realizes the target. Removing either diagnostic would make the inverse claim
less interpretable: surrogate-only optimization hides exploitation, whereas
solver-only reporting hides how the replay budget was allocated.

Several limitations constrain the claim. The inverse experiment uses one
training seed, so its paired significance measures per-target consistency for
that frozen run rather than reproducibility across retraining. All structures
belong to one \(\pgroup\)-symmetric, binary, fixed-material design family on a
single path and resolution. We do not demonstrate transfer to arbitrary
lattices, multiple materials, three dimensions, alternative polarizations, or
fabrication-aware minimum features. The query result also shows poor held-path
extrapolation. Finally, no controlled ablation isolates every component of the
BandBlock; the present evidence supports the implemented BISON system as a
whole, not a causal claim for one gate or spectral branch.

These limitations define concrete next experiments. A component ablation
should remove the second spectral kernel, gate, and low-rank branch one at a
time under the same three-seed forward protocol. A multi-seed inverse study
should retrain both generator variants and repeat the complete 625-target
replay with identical candidate hashes and solver budgets. Resolution and
geometry transfer should compare schema-compatible grids without changing the
reported test identities, while fabrication-aware work should add minimum
feature and robustness constraints before any experimental realization claim.

The dataset is simulated and contains no personal or human-subject data. The
main broader risk is over-trusting a learned surrogate, which could waste
solver, fabrication, or laboratory resources. We mitigate this by publishing
surrogate and replay errors separately, reporting negative extrapolation and
worst-case examples, and requiring physical replay before describing a design
as realized.
